% !TeX root = main.tex
% 五章目錄骨架，對應 ../論文章節架構.md v0.5（2026-11-21）。
% 只有章節標題，未搬入或改寫既有 Markdown 正文。
% 請在 thesis.tex 的 \mainmatter 之後以 % !TeX root = main.tex
% 五章目錄骨架，對應 ../論文章節架構.md v0.5（2026-11-21）。
% 只有章節標題，未搬入或改寫既有 Markdown 正文。
% 請在 thesis.tex 的 \mainmatter 之後以 % !TeX root = main.tex
% 五章目錄骨架，對應 ../論文章節架構.md v0.5（2026-11-21）。
% 只有章節標題，未搬入或改寫既有 Markdown 正文。
% 請在 thesis.tex 的 \mainmatter 之後以 % !TeX root = main.tex
% 五章目錄骨架，對應 ../論文章節架構.md v0.5（2026-11-21）。
% 只有章節標題，未搬入或改寫既有 Markdown 正文。
% 請在 thesis.tex 的 \mainmatter 之後以 \input{chapters_outline} 載入。
% 本檔不包含 \documentclass、\begin{document} 或 \tableofcontents。
% 正式改用各章正文檔時，請移除本檔的載入，避免章節重複。

\chapter{緒論}

\section{研究背景}
\subsection{自主移動機械臂之發展與應用}
\subsection{動態環境下之導航與操作需求}

\section{研究動機與目的}
\subsection{導航與操作整合之挑戰}
\subsection{命令執行差異對預測控制之影響}
\subsection{研究目的與預期貢獻}

\section{文獻回顧}
\subsection{預測導航與動態避障}
\subsection{全身模型預測控制與移動中操作}
\subsection{視覺導引與關節式物件操作}
\subsection{執行動態建模與追蹤補償}
\subsection{現有研究比較與研究缺口}

\section{本文架構}

\chapter{研究系統介紹}

\section{移動機械臂平台}
\subsection{全向移動底盤}
\subsection{六軸機械臂與夾爪}

\section{感測設備與命令介面}
\subsection{雷射掃描與深度相機}
\subsection{底盤里程與關節狀態回授}
\subsection{底盤速度與手臂位置命令介面}

\section{軟體架構與模擬環境}
\subsection{ROS 2 軟體環境與通訊架構}
\subsection{Isaac Sim 物理模擬與場景建置}
\subsection{資料記錄與實驗重現方式}

\section{導航、視覺與操作之系統整合}
\subsection{全域路徑規劃與導航參考}
\subsection{視覺感知與操作目標資訊}
\subsection{單一全身控制器與命令執行架構}

\section{問題定義與基本假設}
\subsection{研究問題與控制目標}
\subsection{環境、感測與已知物件資訊}
\subsection{執行介面與操作模型假設}

\section{任務流程與研究範圍}
\subsection{導航、接近與抓持流程}
\subsection{開門、穿門與取物返回流程}
\subsection{研究範圍與驗證邊界}

\chapter{全身幾何模型預測控制方法}

\section{底盤幾何運動與全身運動學}
\subsection{座標系定義與位姿表示}
\subsection{底盤 SE(2) 運動模型與追蹤誤差}
\subsection{移動機械臂正向運動學與全身雅可比矩陣}

\section{手臂設定點追蹤與增廣預測模型}
\subsection{手臂位置命令與伺服追蹤動態}
\subsection{實際關節位置與設定點之增廣狀態}
\subsection{多步預測模型建立與參數辨識}

\section{全身預測控制之目標函數與限制條件}
\subsection{底盤與末端位姿追蹤目標}
\subsection{關節、命令變化與輪級運動限制}
\subsection{最佳化問題建構與數值求解}
\subsection{滾動時域控制原則}

\section{動態障礙預測與導航參考調整}
\subsection{障礙物表示與運動估測}
\subsection{未來相對運動與衝突預測}
\subsection{速度與側移候選評估及導航參考更新}
\subsection{控制屏障函數與避障限制}
\subsection{原始雷射掃描防護與命令處理}

\section{視覺目標估測與任務參考生成}
\subsection{任務目標之視覺定位與幾何估測}
\subsection{影像、機器人狀態與座標轉換之時間對齊}
\subsection{可行操作區域與底盤目標選擇}
\subsection{導航路徑下之朝向與末端參考生成}

\section{移動中操作與機構約束控制}
\subsection{移動中接近與抓持條件}
\subsection{門與抽屜之機構約束建模}
\subsection{持物運送與收納之運動參考}
\subsection{任務階段切換與控制狀態延續}

\section{執行動態補償與控制實現}
\subsection{在途命令與實際放行輸入之預測整合}
\subsection{模型偏置估測與末端追蹤補償}
\subsection{命令放行、停止軌跡檢查與逾時處置}
\subsection{整體控制流程與演算法}

\chapter{實驗方法與結果分析}

\section{實驗環境與評估指標}
\subsection{實驗平台與場景配置}
\subsection{比較方法與控制條件}
\subsection{誤差、平順性與任務表現指標}

\section{執行動態模型之辨識與預測驗證}
\subsection{手臂追蹤動態參數辨識}
\subsection{理想運動模型與增廣模型之比較}
\subsection{命令延遲與偏置估測之影響}

\section{動態導航與避障比較}
\subsection{動態障礙場景與測試配置}
\subsection{不同導航方法之性能比較}
\subsection{控制屏障與雷射防護之消融比較}

\section{全身協同操作與預測時域比較}
\subsection{底盤與末端協同追蹤表現}
\subsection{單步與多步控制之比較}
\subsection{抽屜操作之移動與停駐策略比較}

\section{住宅場景整合任務}
\subsection{視覺導引之開門與穿門}
\subsection{茶几物件抓取與持物返回}
\subsection{任務銜接與整體完成情形}

\section{計算效能、失敗案例與研究限制}
\subsection{求解時間與控制期限分析}
\subsection{誤差來源與失敗案例分析}
\subsection{方法適用範圍與限制}

\chapter{結論與未來工作}

\section{研究成果與貢獻}
\subsection{幾何預測導航與動態避障}
\subsection{執行動態建模與全身協同控制}
\subsection{視覺導引與受約束操作整合}

\section{未來研究方向}
\subsection{實機模型辨識與即時控制驗證}
\subsection{真實抓持與接觸動態處理}
\subsection{視覺感知與任務泛化能力}

% 參考文獻與附錄待內容備妥再加入；目前不建立空白章頁。
 載入。
% 本檔不包含 \documentclass、\begin{document} 或 \tableofcontents。
% 正式改用各章正文檔時，請移除本檔的載入，避免章節重複。

\chapter{緒論}

\section{研究背景}
\subsection{自主移動機械臂之發展與應用}
\subsection{動態環境下之導航與操作需求}

\section{研究動機與目的}
\subsection{導航與操作整合之挑戰}
\subsection{命令執行差異對預測控制之影響}
\subsection{研究目的與預期貢獻}

\section{文獻回顧}
\subsection{預測導航與動態避障}
\subsection{全身模型預測控制與移動中操作}
\subsection{視覺導引與關節式物件操作}
\subsection{執行動態建模與追蹤補償}
\subsection{現有研究比較與研究缺口}

\section{本文架構}

\chapter{研究系統介紹}

\section{移動機械臂平台}
\subsection{全向移動底盤}
\subsection{六軸機械臂與夾爪}

\section{感測設備與命令介面}
\subsection{雷射掃描與深度相機}
\subsection{底盤里程與關節狀態回授}
\subsection{底盤速度與手臂位置命令介面}

\section{軟體架構與模擬環境}
\subsection{ROS 2 軟體環境與通訊架構}
\subsection{Isaac Sim 物理模擬與場景建置}
\subsection{資料記錄與實驗重現方式}

\section{導航、視覺與操作之系統整合}
\subsection{全域路徑規劃與導航參考}
\subsection{視覺感知與操作目標資訊}
\subsection{單一全身控制器與命令執行架構}

\section{問題定義與基本假設}
\subsection{研究問題與控制目標}
\subsection{環境、感測與已知物件資訊}
\subsection{執行介面與操作模型假設}

\section{任務流程與研究範圍}
\subsection{導航、接近與抓持流程}
\subsection{開門、穿門與取物返回流程}
\subsection{研究範圍與驗證邊界}

\chapter{全身幾何模型預測控制方法}

\section{底盤幾何運動與全身運動學}
\subsection{座標系定義與位姿表示}
\subsection{底盤 SE(2) 運動模型與追蹤誤差}
\subsection{移動機械臂正向運動學與全身雅可比矩陣}

\section{手臂設定點追蹤與增廣預測模型}
\subsection{手臂位置命令與伺服追蹤動態}
\subsection{實際關節位置與設定點之增廣狀態}
\subsection{多步預測模型建立與參數辨識}

\section{全身預測控制之目標函數與限制條件}
\subsection{底盤與末端位姿追蹤目標}
\subsection{關節、命令變化與輪級運動限制}
\subsection{最佳化問題建構與數值求解}
\subsection{滾動時域控制原則}

\section{動態障礙預測與導航參考調整}
\subsection{障礙物表示與運動估測}
\subsection{未來相對運動與衝突預測}
\subsection{速度與側移候選評估及導航參考更新}
\subsection{控制屏障函數與避障限制}
\subsection{原始雷射掃描防護與命令處理}

\section{視覺目標估測與任務參考生成}
\subsection{任務目標之視覺定位與幾何估測}
\subsection{影像、機器人狀態與座標轉換之時間對齊}
\subsection{可行操作區域與底盤目標選擇}
\subsection{導航路徑下之朝向與末端參考生成}

\section{移動中操作與機構約束控制}
\subsection{移動中接近與抓持條件}
\subsection{門與抽屜之機構約束建模}
\subsection{持物運送與收納之運動參考}
\subsection{任務階段切換與控制狀態延續}

\section{執行動態補償與控制實現}
\subsection{在途命令與實際放行輸入之預測整合}
\subsection{模型偏置估測與末端追蹤補償}
\subsection{命令放行、停止軌跡檢查與逾時處置}
\subsection{整體控制流程與演算法}

\chapter{實驗方法與結果分析}

\section{實驗環境與評估指標}
\subsection{實驗平台與場景配置}
\subsection{比較方法與控制條件}
\subsection{誤差、平順性與任務表現指標}

\section{執行動態模型之辨識與預測驗證}
\subsection{手臂追蹤動態參數辨識}
\subsection{理想運動模型與增廣模型之比較}
\subsection{命令延遲與偏置估測之影響}

\section{動態導航與避障比較}
\subsection{動態障礙場景與測試配置}
\subsection{不同導航方法之性能比較}
\subsection{控制屏障與雷射防護之消融比較}

\section{全身協同操作與預測時域比較}
\subsection{底盤與末端協同追蹤表現}
\subsection{單步與多步控制之比較}
\subsection{抽屜操作之移動與停駐策略比較}

\section{住宅場景整合任務}
\subsection{視覺導引之開門與穿門}
\subsection{茶几物件抓取與持物返回}
\subsection{任務銜接與整體完成情形}

\section{計算效能、失敗案例與研究限制}
\subsection{求解時間與控制期限分析}
\subsection{誤差來源與失敗案例分析}
\subsection{方法適用範圍與限制}

\chapter{結論與未來工作}

\section{研究成果與貢獻}
\subsection{幾何預測導航與動態避障}
\subsection{執行動態建模與全身協同控制}
\subsection{視覺導引與受約束操作整合}

\section{未來研究方向}
\subsection{實機模型辨識與即時控制驗證}
\subsection{真實抓持與接觸動態處理}
\subsection{視覺感知與任務泛化能力}

% 參考文獻與附錄待內容備妥再加入；目前不建立空白章頁。
 載入。
% 本檔不包含 \documentclass、\begin{document} 或 \tableofcontents。
% 正式改用各章正文檔時，請移除本檔的載入，避免章節重複。

\chapter{緒論}

\section{研究背景}
\subsection{自主移動機械臂之發展與應用}
\subsection{動態環境下之導航與操作需求}

\section{研究動機與目的}
\subsection{導航與操作整合之挑戰}
\subsection{命令執行差異對預測控制之影響}
\subsection{研究目的與預期貢獻}

\section{文獻回顧}
\subsection{預測導航與動態避障}
\subsection{全身模型預測控制與移動中操作}
\subsection{視覺導引與關節式物件操作}
\subsection{執行動態建模與追蹤補償}
\subsection{現有研究比較與研究缺口}

\section{本文架構}

\chapter{研究系統介紹}

\section{移動機械臂平台}
\subsection{全向移動底盤}
\subsection{六軸機械臂與夾爪}

\section{感測設備與命令介面}
\subsection{雷射掃描與深度相機}
\subsection{底盤里程與關節狀態回授}
\subsection{底盤速度與手臂位置命令介面}

\section{軟體架構與模擬環境}
\subsection{ROS 2 軟體環境與通訊架構}
\subsection{Isaac Sim 物理模擬與場景建置}
\subsection{資料記錄與實驗重現方式}

\section{導航、視覺與操作之系統整合}
\subsection{全域路徑規劃與導航參考}
\subsection{視覺感知與操作目標資訊}
\subsection{單一全身控制器與命令執行架構}

\section{問題定義與基本假設}
\subsection{研究問題與控制目標}
\subsection{環境、感測與已知物件資訊}
\subsection{執行介面與操作模型假設}

\section{任務流程與研究範圍}
\subsection{導航、接近與抓持流程}
\subsection{開門、穿門與取物返回流程}
\subsection{研究範圍與驗證邊界}

\chapter{全身幾何模型預測控制方法}

\section{底盤幾何運動與全身運動學}
\subsection{座標系定義與位姿表示}
\subsection{底盤 SE(2) 運動模型與追蹤誤差}
\subsection{移動機械臂正向運動學與全身雅可比矩陣}

\section{手臂設定點追蹤與增廣預測模型}
\subsection{手臂位置命令與伺服追蹤動態}
\subsection{實際關節位置與設定點之增廣狀態}
\subsection{多步預測模型建立與參數辨識}

\section{全身預測控制之目標函數與限制條件}
\subsection{底盤與末端位姿追蹤目標}
\subsection{關節、命令變化與輪級運動限制}
\subsection{最佳化問題建構與數值求解}
\subsection{滾動時域控制原則}

\section{動態障礙預測與導航參考調整}
\subsection{障礙物表示與運動估測}
\subsection{未來相對運動與衝突預測}
\subsection{速度與側移候選評估及導航參考更新}
\subsection{控制屏障函數與避障限制}
\subsection{原始雷射掃描防護與命令處理}

\section{視覺目標估測與任務參考生成}
\subsection{任務目標之視覺定位與幾何估測}
\subsection{影像、機器人狀態與座標轉換之時間對齊}
\subsection{可行操作區域與底盤目標選擇}
\subsection{導航路徑下之朝向與末端參考生成}

\section{移動中操作與機構約束控制}
\subsection{移動中接近與抓持條件}
\subsection{門與抽屜之機構約束建模}
\subsection{持物運送與收納之運動參考}
\subsection{任務階段切換與控制狀態延續}

\section{執行動態補償與控制實現}
\subsection{在途命令與實際放行輸入之預測整合}
\subsection{模型偏置估測與末端追蹤補償}
\subsection{命令放行、停止軌跡檢查與逾時處置}
\subsection{整體控制流程與演算法}

\chapter{實驗方法與結果分析}

\section{實驗環境與評估指標}
\subsection{實驗平台與場景配置}
\subsection{比較方法與控制條件}
\subsection{誤差、平順性與任務表現指標}

\section{執行動態模型之辨識與預測驗證}
\subsection{手臂追蹤動態參數辨識}
\subsection{理想運動模型與增廣模型之比較}
\subsection{命令延遲與偏置估測之影響}

\section{動態導航與避障比較}
\subsection{動態障礙場景與測試配置}
\subsection{不同導航方法之性能比較}
\subsection{控制屏障與雷射防護之消融比較}

\section{全身協同操作與預測時域比較}
\subsection{底盤與末端協同追蹤表現}
\subsection{單步與多步控制之比較}
\subsection{抽屜操作之移動與停駐策略比較}

\section{住宅場景整合任務}
\subsection{視覺導引之開門與穿門}
\subsection{茶几物件抓取與持物返回}
\subsection{任務銜接與整體完成情形}

\section{計算效能、失敗案例與研究限制}
\subsection{求解時間與控制期限分析}
\subsection{誤差來源與失敗案例分析}
\subsection{方法適用範圍與限制}

\chapter{結論與未來工作}

\section{研究成果與貢獻}
\subsection{幾何預測導航與動態避障}
\subsection{執行動態建模與全身協同控制}
\subsection{視覺導引與受約束操作整合}

\section{未來研究方向}
\subsection{實機模型辨識與即時控制驗證}
\subsection{真實抓持與接觸動態處理}
\subsection{視覺感知與任務泛化能力}

% 參考文獻與附錄待內容備妥再加入；目前不建立空白章頁。
 載入。
% 本檔不包含 \documentclass、\begin{document} 或 \tableofcontents。
% 正式改用各章正文檔時，請移除本檔的載入，避免章節重複。

\chapter{緒論}

\section{研究背景}
\subsection{自主移動機械臂之發展與應用}
\subsection{動態環境下之導航與操作需求}

\section{研究動機與目的}
\subsection{導航與操作整合之挑戰}
\subsection{命令執行差異對預測控制之影響}
\subsection{研究目的與預期貢獻}

\section{文獻回顧}
\subsection{預測導航與動態避障}
\subsection{全身模型預測控制與移動中操作}
\subsection{視覺導引與關節式物件操作}
\subsection{執行動態建模與追蹤補償}
\subsection{現有研究比較與研究缺口}

\section{本文架構}

\chapter{研究系統介紹}

\section{移動機械臂平台}
\subsection{全向移動底盤}
\subsection{六軸機械臂與夾爪}

\section{感測設備與命令介面}
\subsection{雷射掃描與深度相機}
\subsection{底盤里程與關節狀態回授}
\subsection{底盤速度與手臂位置命令介面}

\section{軟體架構與模擬環境}
\subsection{ROS 2 軟體環境與通訊架構}
\subsection{Isaac Sim 物理模擬與場景建置}
\subsection{資料記錄與實驗重現方式}

\section{導航、視覺與操作之系統整合}
\subsection{全域路徑規劃與導航參考}
\subsection{視覺感知與操作目標資訊}
\subsection{單一全身控制器與命令執行架構}

\section{問題定義與基本假設}
\subsection{研究問題與控制目標}
\subsection{環境、感測與已知物件資訊}
\subsection{執行介面與操作模型假設}

\section{任務流程與研究範圍}
\subsection{導航、接近與抓持流程}
\subsection{開門、穿門與取物返回流程}
\subsection{研究範圍與驗證邊界}

\chapter{全身幾何模型預測控制方法}

\section{底盤幾何運動與全身運動學}
\subsection{座標系定義與位姿表示}
\subsection{底盤 SE(2) 運動模型與追蹤誤差}
\subsection{移動機械臂正向運動學與全身雅可比矩陣}

\section{手臂設定點追蹤與增廣預測模型}
\subsection{手臂位置命令與伺服追蹤動態}
\subsection{實際關節位置與設定點之增廣狀態}
\subsection{多步預測模型建立與參數辨識}

\section{全身預測控制之目標函數與限制條件}
\subsection{底盤與末端位姿追蹤目標}
\subsection{關節、命令變化與輪級運動限制}
\subsection{最佳化問題建構與數值求解}
\subsection{滾動時域控制原則}

\section{動態障礙預測與導航參考調整}
\subsection{障礙物表示與運動估測}
\subsection{未來相對運動與衝突預測}
\subsection{速度與側移候選評估及導航參考更新}
\subsection{控制屏障函數與避障限制}
\subsection{原始雷射掃描防護與命令處理}

\section{視覺目標估測與任務參考生成}
\subsection{任務目標之視覺定位與幾何估測}
\subsection{影像、機器人狀態與座標轉換之時間對齊}
\subsection{可行操作區域與底盤目標選擇}
\subsection{導航路徑下之朝向與末端參考生成}

\section{移動中操作與機構約束控制}
\subsection{移動中接近與抓持條件}
\subsection{門與抽屜之機構約束建模}
\subsection{持物運送與收納之運動參考}
\subsection{任務階段切換與控制狀態延續}

\section{執行動態補償與控制實現}
\subsection{在途命令與實際放行輸入之預測整合}
\subsection{模型偏置估測與末端追蹤補償}
\subsection{命令放行、停止軌跡檢查與逾時處置}
\subsection{整體控制流程與演算法}

\chapter{實驗方法與結果分析}

\section{實驗環境與評估指標}
\subsection{實驗平台與場景配置}
\subsection{比較方法與控制條件}
\subsection{誤差、平順性與任務表現指標}

\section{執行動態模型之辨識與預測驗證}
\subsection{手臂追蹤動態參數辨識}
\subsection{理想運動模型與增廣模型之比較}
\subsection{命令延遲與偏置估測之影響}

\section{動態導航與避障比較}
\subsection{動態障礙場景與測試配置}
\subsection{不同導航方法之性能比較}
\subsection{控制屏障與雷射防護之消融比較}

\section{全身協同操作與預測時域比較}
\subsection{底盤與末端協同追蹤表現}
\subsection{單步與多步控制之比較}
\subsection{抽屜操作之移動與停駐策略比較}

\section{住宅場景整合任務}
\subsection{視覺導引之開門與穿門}
\subsection{茶几物件抓取與持物返回}
\subsection{任務銜接與整體完成情形}

\section{計算效能、失敗案例與研究限制}
\subsection{求解時間與控制期限分析}
\subsection{誤差來源與失敗案例分析}
\subsection{方法適用範圍與限制}

\chapter{結論與未來工作}

\section{研究成果與貢獻}
\subsection{幾何預測導航與動態避障}
\subsection{執行動態建模與全身協同控制}
\subsection{視覺導引與受約束操作整合}

\section{未來研究方向}
\subsection{實機模型辨識與即時控制驗證}
\subsection{真實抓持與接觸動態處理}
\subsection{視覺感知與任務泛化能力}

% 參考文獻與附錄待內容備妥再加入；目前不建立空白章頁。
